\documentclass{jsarticle}
\usepackage{amsmath}
\usepackage{amssymb}
\begin{document}
\title{六星占術の運命星と宿命星の求め方}
\author{細木数子、山口雅旭、苫米地英人}
\date{}
\maketitle



  $$ wxyz \cong wx \oplus yz, wx \oplus yz = w'x' \otimes y'z' $$
  $$ x|_{1979} \equiv 19 \oplus 16, 19 \oplus 16 = 35 $$ 
  $$ 19 \oplus 16 = 36 - 1, 36 + 2 = 38, 1979 \mod n = 2 $$
  $$ 38 \otimes 2 = 76, x|_{1940,1979} = 24, x| = {76 - 12 \over 2}, x|_{1979} = 32 \oplus 6 $$
  易学の６４卦を３８周期で求めると、同型写像から３２となり、この真逆が六星占術の暦となる。
  日にちの干支と十干が月の干支と十干となり、真逆が六星占術となる。
  1940と1979の３８周期から干支の素因数分解の剰余項と
  1979の準同形写像からも運命星がもとまる。この運命星は月運の性質でもある。
  $ x = 2n - 2 $は、5行の月運の十干を表している。
  1979年周期で、西暦の加群分解がある。この1979年前は、-1979年で
  0年とおけて、旧約聖書の12月が、0年の2月が36より、1月が5から5-31で-26から
  60-26より12月は34となり、34+31=65によって、0年の1月の生命エネルギーは
  5となり、キリストが生誕した、B.C.4年から、　1年ごとに5と干支と十干が動く
  このから、キリストを1とすると、5年前は0年となり、-1979年と同型から
  -1979年の0年は2月が38となる。1979年周期でこの2倍の2958年の加群分解は、
  29+58=97-60=37となり、この2958年からの2倍の1979年の3倍は、5937=59+37=96-60=39
  となり、神様と人類の即身成仏の数秘術となる。加群分解が成り立つのは、
  この1979を起点とする倍数の±2前後の西暦の数値となっている。
  六星占術は、生命が誕生したら-1加群となる。これが運命数から運命星となる。
  運命から宿命そして立命となる。五行で運気が悪いのが六星占術ではよい運気になっている。
  六星占術で運気が悪いのが五行では運気がよい。易学を細分化すると六星占術となる。
  相性殺界となる運気が六星占術と五行の相殺となる。
  宇宙が始まって1940と1979の38周期が易学の同型写像となり、六星占術となる。
  六星占術の
  0から9までは、土星の陽における種子から陰における安定
  10から19までは、金星の陽の種子から陰の安定
  20から29までは、火星の陽の種子から陰の安定
  30から39までは、天王星の陽の種子から陰の安定
  40から49までは、木星の陽の種子から陰の安定
  50から59までは、水星の陽の種子から陰の安定


  5行では、金星から木星、火星、土星の申から未までを通っている。
  六星占術の大殺界は、本来真逆の財になっている。
  九紫、大黒、緑水星、大善星、妙雅、火竹、白水、静雲、白鴎、光美と
  宿命星が廻っている。
  各宿命星の各星人の周りは別にこの通りではない。　

  $$ x|_{1979} = 35, x = 36 + \sum_{k=1}^{30} x_1 + \sum_{k=1}^{31} x_2 + \sum_{k=1}^{28} x_3 + x_4 $$
  $$ x = 36 + {n (n - 1) \over 2} x_9 + x_4 $$
  $$ x_5 = {x \over 60}, x_6 = {x_5 \over 12}, x_7 = {x_6 \over 10} $$
  $$ x_8 = {2x + x^{'} -1 \over 2} $$

  西暦1979年を西暦0年と生命エネルギー36を閾値としての、各生命エネルギーの
  上下の平行線となっていて、各星人の生命エネルギーは、この線の上と下の
  行き来となる周期で表されている。
  
  六星占術の始まりの年の月における日の求め方は、　
  $$ -1979 \cong 1979, 1979 = 19 \oplus 79 \cong 19 \oplus 16 = 35 \cong 38 $$
  1979年の日の五行の月は、
  $$ x = 2n - 2 $$
  と表せられて、
  2月から宿命星の年が始まることにより、
  求める星人の運命星は、運命数-1加群より、求める生命エネルギーをxとすると、
  $$ x = 35(2月の生命エネルギーから) + \sum(1,3,5,7,8,10,12) \sum(4,6,9,11) 
  + (うるう年の4年に一回+1) $$
  ここで、うるう年が4年に一回なので、xが閏年があるとしだと、2月29日として、
  4年に1回、2月に+1する。
  キリスト誕生の年がB.C.4年で、キリストが旧約聖書のときの西暦0年の前の
  12月25日が-1979年とすると、西暦0年の2月も生命エネルギーが35前後になる。
  神様と人類の即身成仏の数秘術が六星占術となる。

  天運、本人の月運、地運は、六星占術は、地運の誕生のときが生命エネルギー
  としての本人の月運となり、この運命星の十二支が先祖の十二支によって、
  天運となり、
  この基質によって人生の天運のときの相性、月運、地運が動き、
  六星占術の各星人の適職、才能、特技、性格によって、
  運命星の運勢として運命が流れる。大殺界のときは、得手不得手として、
  流れるが、5行を見ると大丈夫になっている。
  各星人の得手、不得手によって、方位が凶か吉として巡る。
  九星は六星占術の相性殺界を使って、凶を吉にしている大人の占いでもある。
  宿命星は、運命数の単体量をだして、この数値に+10として十の宿命星の
  流れを占っている。運命が固定されている。
  要するに、日の各運命数が十干の月運2n-2と同じく、±2の範囲で動き、
  十干と十二支によって決まっているのとは、違う。
  その上に、常占術としての月運が1年に12あるのとは違い、運命星の日の運命数の
  星人と十二支を月に移動させて、先祖の天運の十二支をこの日の星人と干支を
  月とみなしてので、各星人が同じ月なのに違う占いとして成り立っている。








\end{document}
